% UTF-8; XeLaTeX. Keep only packages actually used by this report.
\documentclass[UTF8,a4paper,10pt]{ctexart}
\usepackage[left=3.17cm,right=3.17cm,top=2.74cm,bottom=2.74cm]{geometry}
\usepackage{amsmath,array,graphicx,float,booktabs,xcolor,listings}
\usepackage[colorlinks,linkcolor=black,urlcolor=blue,citecolor=black]{hyperref}
\usepackage{fancyhdr}
\setlength{\headheight}{15pt}
\pagestyle{fancy}\lhead{\kaishu\leftmark}\rhead{\kaishu 编译原理实验报告}\cfoot{\thepage}
\newcommand{\HRule}{\rule{\linewidth}{0.5mm}}
\newcommand{\sihao}{\fontsize{14pt}{21pt}\selectfont}
\ctexset{section={name={,、},number=\chinese{section}},subsection={name={（,）},number=\chinese{subsection}}}
\lstset{basicstyle=\small\ttfamily,breaklines=true,frame=single,columns=fullflexible,keepspaces=true}
\setlength{\emergencystretch}{2em}
\begin{document}
\begin{titlepage}
    \begin{center}
    \includegraphics[width=0.8\textwidth]{NKU.png}\\[1cm]    
    \textsc{\Huge \kaishu{\textbf{南\ \ \ \ \ \ 开\ \ \ \ \ \ 大\ \ \ \ \ \ 学}} }\\[0.9cm]
    \textsc{\huge \kaishu{\textbf{计\ \ 算\ \ 机\ \ 学\ \ 院}}}\\[0.5cm]
    \textsc{\Large \textbf{编译原理上机实验报告}}\\[0.8cm]
    \HRule \\[0.9cm]
    { \LARGE \bfseries 预备工作：了解你的编译器、LLVM IR 与汇编编程}\\[0.4cm]
    \HRule \\[2.0cm]
    \centering
    \textsc{\LARGE \kaishu{姓名\ :\ 张龙飞、袁煜杰}}\\[0.5cm]
    \textsc{\LARGE \kaishu{学号\ :\ 2413625、2412303}}\\[0.5cm]
    \textsc{\LARGE \kaishu{专业\ :\ 计算机科学与技术}}\\[0.5cm]
    \vfill
    {\Large \today}
    \end{center}
\end{titlepage}
%-------------摘------要--------------
\newpage
\thispagestyle{empty}
\renewcommand{\abstractname}{\kaishu \sihao \textbf{摘要}}
    \begin{abstract}
        % 【待补全】摘要正文：概述本次预备工作的目的、内容、方法与结果。
        % 要点提示：以 GCC/LLVM(Clang) 为研究对象，利用命令行选项获得预处理、编译、
        % 汇编、链接各阶段的输出，分析其与源程序的关系；设计覆盖 SysY 语言特性的
        % 示例程序，手写等价的 LLVM IR 与 RISC-V 汇编程序，链接 SysY 运行时库并
        % 在 QEMU 模拟器上验证运行结果，为后续 SysY 编译器的实现奠定基础。
        本报告记录两人的预备实验。袁煜杰完成六组手写 LLVM IR，与参考程序、本机和 RISC-V 执行结果比较，236 项检查全部通过；MLIR 进阶部分保存 Linalg 到 LLVM 的逐层输出，并验证 VecAdd 的 16 个元素。张龙飞已有编译流程与循环阶乘汇编实验，相关分析和更多手写汇编仍待其完善。本版本为共同报告整合稿。

        \noindent  %顶格
        \textbf{\\\ 关键字：编译器；LLVM IR；RISC-V 汇编；SysY；链接}\textbf{} \\\ \\\
    \end{abstract}
%----------------------------------------------------------------
\tableofcontents
%----------------------------------------------------------------
\newpage

\setcounter{page}{1}

%==================================================================
\section{引言}
%==================================================================

\subsection{实验背景与目的}
% 【待补全】用自己的话概述，可包含以下要点：
SysY 是全国大学生编译系统设计赛定义的简化 C 语言。本课程上机大作业要求实现一个
SysY 编译器（词法分析、语法分析、类型检查与中间代码生成、代码优化、目标代码生成），
目标平台为 RISC-V。本次预备工作是上机大作业的第一个模块，目的在于：
\begin{enumerate}
  \item 以 GCC / LLVM(Clang) 为研究对象，深入探究语言处理系统的完整工作过程，
        理解预处理器、编译器（前端与后端各阶段）、汇编器、链接器的功能；
  \item 熟悉 LLVM IR 中间语言与 RISC-V 汇编语言，
        设计覆盖 SysY 语言特性的示例程序，手写等价的 LLVM IR 与 RISC-V 汇编程序，
        链接 SysY 运行时库并验证运行结果的正确性；
  \item 为后续词法分析器、语法分析器、中间代码生成与目标代码生成等模块的实现
        建立直观认识与技术储备。
\end{enumerate}

\subsection{实验环境（张龙飞）}
本实验在 WSL2（Windows Subsystem for Linux）中进行，使用 QEMU 用户态模拟器
运行 RISC-V 目标程序。主要工具及版本如表\ref{tab:env}所示。

\begin{table}[!htbp]
  \centering
  \begin{tabular}{ll}
  \toprule
  项目 & 版本 / 说明\\
  \midrule
  操作系统 & WSL2：Ubuntu 22.04.5 LTS（x86\_64）\\
  GCC / G++ & 11.4.0\\
  Clang / LLVM 工具链 & 14.0.0（含 llvm-as、lli、llc、opt、llvm-dis）\\
  Flex / Bison & 2.6.4 / 3.8.2\\
  RISC-V 交叉工具链 & riscv64-linux-gnu-gcc / g++ 11.4.0\\
  模拟器 & QEMU 7.0.0（qemu-riscv64 用户态模式，/opt/qemu 源码编译）\\
  \bottomrule
  \end{tabular}
  \caption{实验环境与工具版本}
  \label{tab:env}
\end{table}
\subsection{实验环境（袁煜杰）}
实际运行环境为 WSL Ubuntu 26.04.1 LTS（x86\_64），Clang/LLVM 21.1.8，
RISC-V GCC 15.2.0，QEMU 10.2.1。MLIR 21.1.8 工具从 Ubuntu 官方源下载并在项目内解压。
详细命令和版本保存在实验结果目录中。

\subsection{小组分工}
袁煜杰（ashalghy-hub）负责基本要求二的 LLVM IR 手写程序与验证、进阶 MLIR 探索
及相应报告；张龙飞（3769lkj）负责基本要求二的手写 RISC-V 汇编，
并继续完善已有工作流程实验分析。两人共同整合一份报告，开头说明分工，分别提交。
本版本保留搭档尚未补齐的内容，完整的袁煜杰部分另提供独立附件。

\section{编译系统工作流程研究（张龙飞）}
%==================================================================
% 本大题两人各自独立完成、独立撰写。以下骨架基于 PPT 示例的 C++ 版
% 阶乘程序 fact.cpp，命令与数据均已实测，可直接替换为你们自己的截图与观察。

本小题以 PPT 示例的 C++ 版阶乘程序 \texttt{fact.cpp} 为研究载体，利用
编译器的命令行选项逐步获得各阶段输出（实验方案见表\ref{tab:cmd}），
研究其与源程序的关系。

\subsection{实验方案}

\begin{lstlisting}[title={研究载体：阶乘程序 fact.cpp（PPT 示例）},language={C++}]
#include <iostream>
using namespace std;
int main()
{
    int i, n, f;

    cin >> n;
    i = 2;
    f = 1;
    while (i <= n)
    {
        f = f * i;
        i = i + 1;
    }
    cout << f << endl;
}
\end{lstlisting}

\begin{table}[!htbp]
  \centering
  \begin{tabular}{lll}
  \toprule
  阶段 & 命令选项 & 产物\\
  \midrule
  预处理 & \texttt{clang++ -E} & \texttt{fact.i}\\
  编译（前端） & \texttt{clang++ -emit-llvm -S} & \texttt{fact.ll}（LLVM IR）\\
  编译（后端） & \texttt{g++ -S} / \texttt{riscv64-linux-gnu-g++ -S} & \texttt{fact.s}（汇编）\\
  汇编 & \texttt{clang++ -c} & \texttt{fact.o}（目标文件）\\
  链接 & \texttt{g++} / \texttt{riscv64-linux-gnu-g++ -static} & 可执行程序\\
  运行验证 & \texttt{lli} / \texttt{qemu-riscv64} & 运行结果\\
  \bottomrule
  \end{tabular}
  \caption{各编译阶段对应的命令选项与产物}
  \label{tab:cmd}
\end{table}

\subsection{预处理器}
% 【待补全】你们的观察与分析。要点提示：
% 宏展开、头文件递归展开（stdio.h 注入大量声明）、注释删除、条件编译处理。
源程序仅 17 行，经 \texttt{clang++ -E} 预处理后得到 30194 行的
\texttt{fact.i}——C++ 标准库头文件 \texttt{<iostream>} 展开了数万行声明。

\begin{figure}[H]
    \centering
    \fbox{\parbox[c][4cm][c]{0.85\textwidth}{\centering 【待插入截图：fact.i 与 fact.cpp 的内容对比】}}
    \caption{预处理输出与源程序的对比}
    \label{fig:prep}
\end{figure}

【待补全：预处理前后对比分析，例如 \texttt{\#include} 与 \texttt{\#define}
的处理、注释的去留、可用 \texttt{-D} 选项自定义宏做进一步观察，以及
C++ 头文件展开规模远大于 C 头文件的原因】

\subsection{编译器}
\subsubsection{编译前端：生成 LLVM IR}
% 【待补全】要点提示：alloca/load/store 与局部变量的对应；icmp sle + br 与
% while 循环的对应；cin/cout 变为 std::istream/ostream 的成员函数调用
% （C++ 符号修饰，如 _ZNSirsERi）；全局量 @_ZSt3cin/@_ZSt4cout 与
% __cxx_global_var_init 初始化函数；类型系统 i32 等。
\texttt{clang++ -emit-llvm -S} 生成 111 行 LLVM IR（\texttt{fact.ll}），
关键片段如下：

\begin{lstlisting}[title={fact.ll 关键片段}]
@_ZSt3cin  = external global %"class.std::basic_istream", align 8
@_ZSt4cout = external global %"class.std::basic_ostream", align 8

define dso_local noundef i32 @main() #4 {
  %1 = alloca i32, align 4      ; return value
  %2 = alloca i32, align 4      ; i
  %3 = alloca i32, align 4      ; n
  %4 = alloca i32, align 4      ; f
  store i32 0, i32* %1, align 4
  %5 = call ... @_ZNSirsERi(@_ZSt3cin, i32* %3)      ; cin >> n
  store i32 2, i32* %2, align 4                      ; i = 2
  store i32 1, i32* %4, align 4                      ; f = 1
  br label %6
6:                                                   ; while (i <= n)
  %7 = load i32, i32* %2, align 4
  %8 = load i32, i32* %3, align 4
  %9 = icmp sle i32 %7, %8
  br i1 %9, label %10, label %16
10:
  %11 = load i32, i32* %4
  %12 = load i32, i32* %2
  %13 = mul nsw i32 %11, %12                         ; f = f * i
  ...
16:
  %17 = load i32, i32* %4, align 4
  %18 = call ... @_ZNSolsEi(@_ZSt4cout, i32 %17)     ; cout << f
  %19 = call ... @_ZNSolsEPFRSoS_E(%18, ... endl)    ; << endl
  ret i32 %20
}
\end{lstlisting}

\begin{figure}[H]
    \centering
    \fbox{\parbox[c][4cm][c]{0.85\textwidth}{\centering 【待插入：fact.ll 控制流图（已生成：prep/stages/fact/fact\_main\_cfg.png）】}}
    \caption{LLVM IR 与控制流结构}
    \label{fig:ll}
\end{figure}

【待补全：IR 与源程序的对应关系分析】

\subsubsection{编译后端：生成目标汇编}
% 【待补全】要点提示：栈帧建立/销毁；参数按调用约定放入寄存器；
% C++ 成员函数调用经链接器解析（_ZNSirsERi@plt 等）；
% x86(CISC) 与 RISC-V(RISC) 的风格差异（如访存指令、立即数处理）。
分别用 \texttt{g++ -S}（x86\_64，AT\&T 风格）与
\texttt{riscv64-linux-gnu-g++ -S}（RISC-V）生成汇编，对比两者风格。

\begin{figure}[H]
    \centering
    \fbox{\parbox[c][4cm][c]{0.85\textwidth}{\centering 【待插入截图：x86 汇编与 RISC-V 汇编对比】}}
    \caption{同一代码在 x86 与 RISC-V 下的汇编对比}
    \label{fig:asm}
\end{figure}

【待补全：x86 与 RISC-V 汇编的对比分析】

\subsection{汇编器}
% 【待补全】要点提示：汇编指令翻译为机器码；符号表与重定位信息
% （objdump -r 可见 call 的目标地址尚为占位）；file 命令查看 ELF 格式。
\texttt{clang++ -c} 将汇编翻译为目标文件 \texttt{fact.o}，可用
\texttt{objdump -d} 反汇编查看机器码，\texttt{objdump -r} 查看重定位项。

\begin{figure}[H]
    \centering
    \fbox{\parbox[c][3.5cm][c]{0.85\textwidth}{\centering 【待插入截图：objdump -d / objdump -r 输出】}}
    \caption{目标文件的反汇编与重定位信息}
    \label{fig:objdump}
\end{figure}

【待补全：汇编器功能分析】

\subsection{链接器}
% 【待补全】要点提示：std::cin/std::cout 等 C++ 符号在链接时解析到
% libstdc++；动态链接（ldd 可见依赖）vs 静态链接（文件体积大但可独立运行，
% 交叉编译产物用 qemu-riscv64 直接运行）。
链接阶段将目标文件与 C++ 运行库（\texttt{std::cin}/\texttt{std::cout}
所在的 libstdc++）组装为可执行程序。动态链接与静态链接产物对比如
图\ref{fig:link}所示。

\begin{figure}[H]
    \centering
    \fbox{\parbox[c][3.5cm][c]{0.85\textwidth}{\centering 【待插入截图：ldd 输出 / 动态与静态可执行文件大小对比】}}
    \caption{动态链接与静态链接对比}
    \label{fig:link}
\end{figure}

【待补全：链接器功能分析】

\subsection{运行验证}
四种方式运行同一程序，结果一致（输入 5，输出 120），说明各阶段产物等价：

\begin{lstlisting}[title={运行验证},language=bash]
$ echo 5 | ./fact_x86                 # x86 可执行（g++ 链接产物）
120
$ echo 5 | lli fact.ll                # LLVM IR 解释执行
120
$ echo 5 | ./fact_fromll_x86          # 由 IR 直接重编译的可执行
120
$ echo 5 | qemu-riscv64 ./fact_rv64   # RISC-V 可执行（QEMU 模拟）
120
\end{lstlisting}

\subsection{进一步探索：优化选项的影响}
% PPT 期望的加分探索。数据来自 opt_compare/SUMMARY.txt（已实测）。
调整优化级别（\texttt{-O0/-O1/-O2/-Os}），观察 LLVM IR 与 RISC-V 汇编的
变化，结果如表\ref{tab:opt}所示。

\begin{table}[!htbp]
  \centering
  \begin{tabular}{lcccc}
  \toprule
  程序 / 产物 & -O0 & -O1 & -O2 & -Os\\
  \midrule
  fact：LLVM IR 行数 & 111 & 163 & 250 & 152\\
  fact：RISC-V 汇编行数 & 156 & 97 & 146 & 89\\
  fib：LLVM IR 行数 & 124 & 248 & 247 & 153\\
  fib：RISC-V 汇编行数 & 177 & 156 & 231 & 118\\
  \bottomrule
  \end{tabular}
  \caption{不同优化级别下 LLVM IR 与 RISC-V 汇编的规模对比}
  \label{tab:opt}
\end{table}

% 【待补全】分析提示：fact 的 IR 在优化后行数显著增加（内联展开 C++ 库调用、
% 循环展开、向量化等），而 RISC-V 汇编行数下降（冗余 load/store 消除、
% mem2reg 等）；还可对比 g++ 与 clang++ 的 x86 汇编差异（133 行 vs 113 行）、
% 加入 -g 调试选项后 objdump -S 的源码级对照、微调源程序后的输出变化等。
【待补全：结合表\ref{tab:opt}及更多探索（如 \texttt{-g} 选项、微调源程序）
分析优化器行为】

\section{LLVM IR 与汇编编程实践}
\subsection{实验方法}
设计六组 SysY 程序，分别编写 LLVM IR 与 RISC-V 汇编。
整数和浮点数分别用 i32、float 表示；局部寄存器遵循 SSA，控制流合流使用 phi，
内存变量通过 load/store 访问\cite{llvmref}。

\subsection{编译与验证流程}
使用 Clang/LLVM 检查、优化并编译 IR，链接 getint、putint、putch 完成输入输出。
验证脚本比较参考程序、手写 IR 和手写汇编的标准输出及退出码。
\par\vspace{0.7\baselineskip}\noindent\begin{minipage}[t]{\linewidth}
\begin{lstlisting}[language=bash]
python3 prep/llvm_ir/verify.py
llvm-as sample.ll -o sample.bc
opt -passes=verify -disable-output sample.ll
opt -S '-passes=default<O2>' sample.ll -o sample_O2.ll
llc -mtriple=riscv64-linux-gnu -mattr=+m,+a,+f,+d,+c \
    -target-abi=lp64d sample.ll -o sample_generated_rv.s
\end{lstlisting}
\end{minipage}\par\vspace{0.5\baselineskip}

\subsection{示例设计与特性覆盖}
\begin{table}[htbp]
\centering\small
\begin{tabular}{p{0.23\linewidth}p{0.54\linewidth}r}
\toprule
示例 & 主要覆盖 & 输入数\\
\midrule
arithmetic & 加减乘除与余数、单目正负、常量、块遮蔽、逻辑非转整数 & 6\\
short\_circuit & if/else、逻辑与或短路、全局副作用、自定义函数 & 6\\
loop\_control & while、continue、break、合流与多条回边 & 6\\
functions & 带参函数、递归、整数返回、void 函数及返回 & 5\\
array\_float & 二维数组、部分初始化、数组形参、浮点运算与转换 & 5\\
promote\_demo & 内存形式循环及 mem2reg 对照 & 5\\
\bottomrule
\end{tabular}
\caption{本次手写 LLVM IR 的实验设计}
\end{table}

算术例子输入 $n$，计算 $(n+3)\times2-n/2+n\%3+7$，并分别输出
$-n$、$+n$、$n+!n$。内层同名变量固定为 7，外层 $n$ 在后续输出中仍保持原值，
用于验证遮蔽而非覆盖。负数整除按向零截断计算，余数为
$n-\operatorname{trunc}(n/3)\times3$。
逻辑非先由 icmp 产生 i1，再由 zext 转为 i32 后参与加法。

\subsection{短路控制流与可观测副作用}
短路例子包含 \texttt{n != 0 \&\& 12 / n > 2} 和
\texttt{n == 0 || touch()}。touch 每次调用将全局 hits 加一。
以 n=0 为例，正确程序既不除以零，也不调用 touch，最终 hits=0；
非零输入则调用 touch 一次。

\par\vspace{0.7\baselineskip}\noindent\begin{minipage}[t]{\linewidth}
\begin{lstlisting}
%nonzero = icmp ne i32 %n, 0
br i1 %nonzero, label %and.rhs, label %and.false
and.rhs:
  %q = sdiv i32 12, %n
  %large = icmp sgt i32 %q, 2
  br i1 %large, label %and.true, label %and.false
and.join:
  %and.value = phi i32 [1, %and.true], [0, %and.false]
\end{lstlisting}
\end{minipage}\par\vspace{0.5\baselineskip}
右操作数置于独立基本块中，由分支决定是否执行。
测试 n=-3、0、1、4、5、12，覆盖受保护分母、条件真/假和副作用次数。

\subsection{循环、函数与数据表示}
循环例子逐个处理 1 到 n，跳过 2，在数值超过 5 时退出。
continue 的边直接回到循环头，break 的边直接进入出口。
循环头分别对 i 和 sum 建立 phi，continue 边保持 sum，普通边传入累加后的值。
出口的 phi 对正常条件退出和 break 退出分别提供值。

\par\vspace{0.7\baselineskip}\noindent\begin{minipage}[t]{\linewidth}
\begin{lstlisting}
header:
  %i = phi i32 [0, %entry], [%next, %skip], [%next, %accumulate]
  %sum = phi i32 [0, %entry], [%sum, %skip], [%added, %accumulate]
  %more = icmp slt i32 %i, %n
  br i1 %more, label %body, label %done
done:
  %result = phi i32 [%sum, %header], [%sum, %break.test]
\end{lstlisting}
\end{minipage}\par\vspace{0.5\baselineskip}
LLVM verifier 检查类型、前驱和支配关系，运行测试检查输出。输入 n=2 时只累加 1，n=5 或 n=8 时结果均为 13，
分别触发 continue 与 break。函数例子进一步使用递归阶乘和 void 输出函数，
检查参数、call 和 ret 的对应关系。

数组/浮点例子将 \texttt{int a[2][3]={{1,2,3},{4}}} 表示为二维数组对象，
初始化第二行剩余元素为零。数组实参传地址，按 \texttt{[3 x i32]} 的行类型
通过 getelementptr 寻址，两行和为 10。
浮点过程为 $\operatorname{twice}(1.5n+0.5)=3n+1$，使用 sitofp、fmul、fadd、
fptosi 和 fcmp，既验证 int/float 转换，也验证 float 参数和返回值。


\subsection{mem2reg 观察与结果验证}
另手写一份以 alloca/load/store 保存循环变量的求和程序，使用 LLVM 自带
\texttt{opt -passes=mem2reg} 得到可提升变量的寄存器形式。
转换前为 2 个 alloca、3 个 load、4 个 store、0 个 phi；转换后
这些局部访存操作均消失，形成 2 个 phi。对 n=-3、0、1、5、10
分别运行，转换前后输出一致。


\subsection{手写 RISC-V 汇编程序及验证}
% 【小题二任务二负责人撰写】
手写 RISC-V（RV64）汇编遵循标准调用约定（a0--a7 传参、ra 保存返回地址、
s0--s11 为被调者保存寄存器）：

\begin{lstlisting}[title={手写 RISC-V 汇编：fact\_manual\_rv.s（关键片段）}]
main:
    addi  sp, sp, -32
    sd    ra, 24(sp)
    sd    s0, 16(sp)
    sd    s1, 8(sp)
    sd    s2, 0(sp)
    call  getint          # a0 = n
    mv    s0, a0
    li    s1, 2           # i = 2
    li    s2, 1           # f = 1
.Lloop:
    bgt   s1, s0, .Ldone  # while (i <= n)
    mul   s2, s2, s1      # f = f * i
    addi  s1, s1, 1       # i = i + 1
    j     .Lloop
.Ldone:
    mv    a0, s2
    call  putint
    li    a0, 10
    call  putch
    li    a0, 0
    ld    ra, 24(sp)
    ld    s0, 16(sp)
    ld    s1, 8(sp)
    ld    s2, 0(sp)
    addi  sp, sp, 32
    ret
\end{lstlisting}

交叉编译并通过 QEMU 运行，输入 7 输出 5040，结果正确：

\begin{lstlisting}[language=bash]
$ riscv64-linux-gnu-gcc -static fact_manual_rv.s sylib.c -o fact_manual_rv
$ echo 7 | qemu-riscv64 ./fact_manual_rv
5040
\end{lstlisting}

【待补全：更多示例程序的手写汇编与验证过程，可对比手写汇编与
\texttt{riscv64-linux-gnu-gcc -S} 自动生成汇编的差异】

\subsection{共同验收状态}
已有循环阶乘手写汇编完成入门验证。扩展 LLVM IR 的 236 项检查见前述章节，
其中自动生成的 RISC-V 汇编不计入搭档的手写任务。
更多手写汇编需由张龙飞完成后按同源程序、同输入和同预期结果接入验收。

\section{进阶要求：MLIR 多级中间表示探索}
\subsection{研究方法}
依据 AscendNPU IR 官方 VecAdd 示例和方言文档分析中端层级，
并使用 MLIR 21.1.8 构造上游 VecAdd，执行逐层降低和 CPU 运行验证。

\subsection{AscendNPU IR VecAdd 的方言与层级}
官方快速入门\cite{ascendexample}的 VecAdd 使用 i16 元素，并区分 GM 与 UB
地址空间。输入先搬入片上缓冲，再进行向量加法，最后搬回输出；设备入口用
HACC 属性表达。编译后还需要 CANN runtime 注册和启动设备程序。


架构说明\cite{ascendarchitecture}把 HFusion、HIVM、HACC 分为高层融合、
硬件相关计算/搬运/同步、异构执行表达等职责。结合 HIVM 文档\cite{hivmref}，
可以看到片上内存和计算流水信息在特定层被显式保留。
官方快速入门从 HIVM 级程序开始，体现片上缓冲、向量计算和数据搬运的配合。

\begin{table}[htbp]
\centering\small
\begin{tabular}{p{0.22\linewidth}p{0.31\linewidth}p{0.36\linewidth}}
\toprule
比较点 & 上游 CPU 实验 & Ascend 官方文档路线\\
\midrule
高层计算 & Linalg 通用逐元素计算 & HFusion 等较高层表达\\
硬件相关操作 & 普通 MemRef 访存 & HIVM 中的计算、搬运、同步\\
内存空间 & 本机普通内存 & GM、UB 等硬件空间\\
执行接口 & MLIR 生成 C 包装函数 & HACC 入口及 CANN runtime\\
证据 & 本机各层输出与运行检查 & 官方示例和架构文档分析\\
\bottomrule
\end{tabular}
\caption{两种路线的职责与证据差异}
\end{table}

\subsection{可执行的上游 VecAdd 设计}
独立设计的 \texttt{prep/mlir/vecadd.mlir} 使用三个
\texttt{memref<16xi16>} 参数。Linalg 的三个恒等索引映射表示
$c[i]=a[i]+b[i]$，迭代属性 parallel 表明各元素计算互相独立。
本实验将其降低为标量循环。
\par\vspace{0.7\baselineskip}\noindent\begin{minipage}[t]{\linewidth}
\begin{lstlisting}
linalg.generic {
  indexing_maps = [affine_map<(i)->(i)>,
                  affine_map<(i)->(i)>, affine_map<(i)->(i)>],
  iterator_types = ["parallel"]
} ins(%a, %b : memref<16xi16>, memref<16xi16>)
  outs(%c : memref<16xi16>) {
^bb0(%x: i16, %y: i16, %old: i16):
  %sum = arith.addi %x, %y : i16
  linalg.yield %sum : i16
}
\end{lstlisting}
\end{minipage}\par\vspace{0.5\baselineskip}
输入数组设为 $a[i]=i$、$b[i]=1$（$0\leq i<16$），C 驱动逐项检查输出。

\subsection{逐层 lowering 与检查点}
上游 MLIR 文档介绍结构化 SCF 与非结构化 CF 的区分\cite{scfref}。
本次根据工具可用 pass\cite{mlirpasses}执行如下管线，每一步均保存文件：
\par\vspace{0.7\baselineskip}\noindent\begin{minipage}[t]{\linewidth}
\begin{lstlisting}[language=bash]
mlir-opt vecadd.mlir -o 00_linalg.mlir
mlir-opt 00_linalg.mlir --convert-linalg-to-loops -o 01_scf.mlir
mlir-opt 01_scf.mlir --convert-scf-to-cf -o 02_cf.mlir
mlir-opt 02_cf.mlir --convert-to-llvm \
    --reconcile-unrealized-casts -o 03_llvm.mlir
mlir-translate --mlir-to-llvmir 03_llvm.mlir -o 04_llvm.ll
\end{lstlisting}
\end{minipage}\par\vspace{0.5\baselineskip}
各阶段输出由验证脚本保存。

\begin{table}[htbp]
\centering\small
\begin{tabular}{lrl}
\toprule
检查点 & 行数 & 实际观察到的操作方言\\
\midrule
00\_linalg & 12 & Linalg、Arith、Func\\
01\_scf & 15 & SCF、MemRef、Arith、Func\\
02\_cf & 21 & CF、MemRef、Arith、Func\\
03\_llvm & 67 & LLVM\\
\bottomrule
\end{tabular}
\caption{MLIR 降低阶段与操作方言}
\end{table}

首先 generic 被展开为 \texttt{scf.for}，循环体出现两个 memref.load、一次
arith.addi 和一次 memref.store。接着循环被拆成比较、条件分支、循环体和回边，
迭代变量通过块参数传递；结构化循环这个概念不再显式保留。
再向 LLVM 方言转换时，MemRef 描述符中的已分配地址、对齐地址、偏移、长度和步长
被显式展开，生成与 C 调用者衔接的包装函数。
最后 \texttt{mlir-translate} 才输出真正的 LLVM IR。
因此 \texttt{llvm.*} 操作所在的 MLIR 文件，与 LLVM 自身的 \texttt{.ll} 文件
并不是同一种文本格式。

\subsection{末端验证与分析}
C 驱动构造与一维 MemRef 对应的描述符，通过
\texttt{\_mlir\_ciface\_vecadd} 调用降低后的程序。
本次 LLVM verifier 检查通过，Clang 编译链接后实际输出：
\par\vspace{0.7\baselineskip}\noindent\begin{minipage}[t]{\linewidth}
\begin{lstlisting}
1 2 3 4 5 6 7 8 9 10 11 12 13 14 15 16
\end{lstlisting}
\end{minipage}\par\vspace{0.5\baselineskip}
驱动逐个断言全部 16 个元素正确，退出码 0。
版本、命令、stderr 和结果保存于 \path{prep/mlir/results/}。
行数从 12 增至 67，主要来自内存描述符与调用接口的展开。

多级 IR 在高层保留索引映射和计算结构，在中层明确循环与访存，
最终转换为指针操作和调用接口。各阶段的类型、控制流和内存表示随降低同步转换。


\section{总结与体会}
袁煜杰负责的手写 LLVM IR 与 MLIR 实验已取得可复现结果。
SSA 合流、短路求值、数组布局和目标 ABI 需要同时满足结构与语义要求，
多级 IR 降低则展示抽象信息如何逐层转为循环、访存与指针操作。
共同报告仍需搭档补齐工作流程分析及扩展手写汇编，再统一提交。
本机仅使用最小教学运行库，取得课程正式库后还需复测。


\begin{thebibliography}{9}
\bibitem{llvmref} LLVM Project. LLVM Language Reference Manual. \url{https://llvm.org/docs/LangRef.html}.
\bibitem{ascendexample} AscendNPU IR. 快速入门示例. \url{https://ascendnpu-ir.gitcode.com/zh_cn/sources/introduction/quick_start/examples_zh.html}.
\bibitem{ascendarchitecture} AscendNPU IR. 架构介绍. \url{https://ascendnpu-ir.gitcode.com/zh_cn/sources/introduction/architecture_zh.html}.
\bibitem{hivmref} AscendNPU IR. HIVM Dialect. \url{https://ascendnpu-ir.gitcode.com/zh_cn/sources/developer_guide/dialects/HIVMDialect.html}.
\bibitem{scfref} MLIR Project. SCF Dialect. \url{https://mlir.llvm.org/docs/Dialects/SCFDialect/}.
\bibitem{mlirpasses} MLIR Project. Passes. \url{https://mlir.llvm.org/docs/Passes/}.
\end{thebibliography}

\end{document}
